\chapter{Angular Standing Modes and the Origin of Atomic Structure}
\label{app:atomic_angular_modes}
\addtocontents{toc}{\protect\setcounter{tocdepth}{0}}

\noindent\textbf{Abstract:} The TFT Hydrogenic Shell Law (\hyperref[app:hydrogenic_shell]{Appendix \ref*{app:hydrogenic_shell}}) establishes the radial structure of the atom, and the Topological Shell Theorem (Section 5.4) proves the emergence of internal cancellation and external neutrality from a time-averaged spherical shell. Standard quantum mechanics introduces orbital and magnetic quantum numbers $l,m$ through probabilistic angular eigenfunctions. Tension Field Theory (TFT) replaces this abstract construction with a deterministic mechanical interpretation: angular structure arises from the allowed standing-precession harmonics of a localized electron knot moving in a central chiral strain field. This appendix derives the TFT angular standing modes, identifying their quantization conditions as the harmonic decomposition of a localized topological orbit in a contiguous lattice.

\section{Physical Premise}
From Part I, matter is a localized topological resonant knot in the Soft Floor. From Part II, the electron is not a diffuse wave but a localized, phase-locked tethered knot whose $4\pi$ topology gives the mechanical origin of spin-1/2.

A stable atomic shell forms when the circulating electron knot closes its phase after an integer number of orbital action cycles, creating a time-averaged standing-strain envelope. If the radial shell is fixed at $r_n$, the allowed angular structure is the set of mechanically allowed standing-precession modes of the localized electron's strain wake.

\section{Decomposition of Atomic Motion}
The localized electron node in hydrogenic motion has three simultaneous components: radial closure, azimuthal circulation, and orbital-plane precession. TFT replaces the standard $n, l, m$ hierarchy with a direct mechanical decomposition:
\begin{itemize}
    \item $n \leftrightarrow$ radial phase closure
    \item $l \leftrightarrow$ angular standing-precession order
    \item $m \leftrightarrow$ projected orientation index
\end{itemize}

\section{The Angular Strain Field on a Fixed Shell}
Because the shell radius is fixed at $r_n$, the angular strain amplitude $\Psi_n (\theta,\phi,t)$ evolves only on the sphere $S^2$. The natural mechanical equation for a standing mode on a closed spherical shell is:
\begin{equation}
    \frac{\partial^2 \Psi_n}{\partial t^2} = c_{\mathrm{ang}}^2 \nabla_{S^2}^2 \Psi_n
\end{equation}
Where $c_{\mathrm{ang}}$ is the effective angular phase speed along the shell, and $\nabla_{S^2}^2$ is the Laplace--Beltrami operator on the sphere. Regularity on the closed, contiguous sphere requires discrete eigenvalues:
\begin{equation}
    \lambda = l(l+1), \quad l=0,1,2,\dots
\end{equation}
Thus, $l$ is strictly the harmonic order of the standing-precession mode on the shell.

\section{Why the Angular Spectrum Is Discrete}
In TFT, discreteness follows mechanically because the localized electron wake must reinforce itself after each full precessional sweep without tearing the lattice discontinuously.

A non-integer angular mode would fail to close smoothly on the sphere, generating a strain mismatch. Because the Soft Floor is contiguous, such a mismatch cannot persist. The shell admits only integer standing harmonics.

\section{The Mechanical Bound ($l \leq n-1$)}
The shell index $n$ fixes the total phase-action budget available on the shell ($m_{\mathrm{e}} v_n r_n = n\hbar$). Each additional angular node line on the shell represents a real reversal of the strain pattern, consuming a finite amount of this phase budget.

Requiring one quantum of action to maintain radial shell closure, the remaining available angular complexity is bounded:
\begin{equation}
    l_{\mathrm{max}} = n - 1
\end{equation}
This finite-capacity constraint geometrically recovers the standard bound: $l = 0, 1, 2, \dots, n-1$.

\section{The Magnetic Submodes ($m$)}
The standing angular mode may carry an azimuthal winding around a preferred observation axis $z$. Requiring single-valued closure in the azimuthal spatial direction ($\Psi(\theta,\phi+2\pi) = \Psi(\theta,\phi)$) implies an azimuthal factor of $e^{im\phi}$ where $m \in \mathbb{Z}$.

Regularity of the full spherical mode strictly requires:
\begin{equation}
    |m| \leq l
\end{equation}
Therefore, $m$ counts the allowed projected azimuthal windings of the shell strain about the chosen axis. This yields the standard $2l+1$ degeneracy count.

\section{Effective Angular Momentum and Spin Separation}
Because the shell supports angular harmonics labeled by $l$, the effective orbital angular momentum magnitude is:
\begin{equation}
    L^2 = l(l+1)\hbar^2, \quad L_z = m\hbar
\end{equation}
In TFT, these are not intrinsic abstract operators; they are the macroscopic geometric actions associated with the standing-precession mode on the shell.

Crucially, the electron itself already carries intrinsic tethered spin-1/2 ($s$) generated by the $4\pi$ topology of the localized knot. Orbital angular structure belongs to the time-averaged shell strain, while intrinsic spin belongs to the localized electron node itself.
